% Source: https://artofproblemsolving.com/community/c3h1710138_junior_regional_olympiad__fbh_2017_grade_7_problem_2
% Stable identity: aops-topic-1710138-q01
Extend $AD$ and $BC$ to meet at $P$. Triangle $ABP$ has two $60^\circ$ angles, so it is equilateral with side $6$, and its area is $\dfrac{\sqrt3}{4}\cdot36=9\sqrt3$.

$D$ lies on $AP$ with $DP=6-4=2$. Since $\angle ADC=90^\circ$, also $\angle PDC=90^\circ$, so triangle $PDC$ is a $30^\circ$-$60^\circ$-$90^\circ$ triangle with the $60^\circ$ angle at $P$: $PC=2PD=4$ and $DC=2\sqrt3$. Hence $C$ lies on segment $PB$ with $BC=6-4=2$.

\begin{asy}
usepackage("amsmath");
size(6cm);
defaultpen(fontsize(10pt));
real s=sqrt(3);
pair A=(0,0),B=(6,0),C=(5,s),D=(2,2s),P=(3,3s);
draw(D--P--C,dashed+gray);
draw(A--B--C--D--cycle,linewidth(0.9));
draw(A--C,blue);draw(B--D,blue);
pair u=unit(A-D),v=unit(C-D);real q=0.35;
draw(D+q*u--D+q*(u+v)--D+q*v);
draw(arc(A,0.8,0,60));draw(arc(B,0.8,120,180));
label("$60^\circ$",A+1.4*dir(40));label("$60^\circ$",B+1.5*dir(160));
dot(P);
label("$A$",A,SW);label("$B$",B,SE);label("$C$",C,E);label("$D$",D,W);label("$P$",P,N);
label("$6$",(3,0),S);label("$4$",(1,s),NW);label("$2$",(5.5,s/2),E);label("$2$",(2.5,2.5s),W);label("$4$",(4,2s),E);
label("\mbox{}",(min(currentpicture,user=true)+max(currentpicture,user=true))/2,black); // keeps the text after the figure black when LaTeX typesets the labels inline
\end{asy}

\textbf{Diagonals.} In the right triangle $ADC$, $AC^2=AD^2+DC^2=16+12=28$. In triangle $ABD$, by the law of cosines, $BD^2=36+16-2\cdot6\cdot4\cdot\dfrac{1}{2}=28$. So
\[
AC=BD=\boxed{2\sqrt7}.
\]

\textbf{Area.} $[ABCD]=[ABP]-[PDC]=9\sqrt3-\dfrac{1}{2}\cdot2\cdot2\sqrt3=\boxed{7\sqrt3}$.
